%-------------------------------------------------------------------------------
%	SECTION TITLE
%-------------------------------------------------------------------------------
\cvsection{Expérience professionnelle}


%-------------------------------------------------------------------------------
%	CONTENT
%-------------------------------------------------------------------------------
\begin{cventries}

\cventry
    {Chercheur postdoctoral - Modèles vision-langage pour la biométrie} % Job title
    {Institut de recherche Idiap} % Organization
    {Martigny, Suisse} % Location
    {Avr. 2025 - Mars 2026} % Date(s)
    {
      \begin{cvitems} % Description(s) of tasks/responsibilities
        \item {Développement de pipelines d'évaluation de bout en bout pour des modèles vision-langage avec les outils de l'écosystème ML moderne}
        \item {Conception de nouvelles métriques d'évaluation et d'outils d'analyse automatisés pour quantifier les performances des modèles selon plusieurs dimensions de qualité, motivées par les exigences de professionnels légaux}
        \item {Mise en oeuvre de workflows ML reproductibles et de pratiques MLOps : gestion de la configuration (Hydra), suivi des expériences (MLflow), orchestration des pipelines (Snakemake)}
        \item {1\textsuperscript{er} prix (15\,000 CHF) lors d'un hackathon startup de 9 jours en équipe de 3 ; développement du prototype fonctionnel de bout en bout et co-présentation du pitch devant un jury technique et commercial, démontrant une capacité de prototypage rapide et de communication d'une vision technique}
        \item {Coordination d'un projet interdisciplinaire sur une plateforme de collecte de données forensiques, en servant d'interface entre une chercheuse en neurosciences cognitives, une développeuse logiciel et une consultante externe en forces de l'ordre, et en traduisant les besoins métier en spécifications techniques}
      \end{cvitems}
    }

%---------------------------------------------------------
  \cventry
    {Doctorant - Biométrie, Analyse forensique d'images} % Job title
    {Institut de recherche Idiap} % Organization
    {Martigny, Suisse} % Location
    {Jan. 2021 - Mars 2025} % Date(s)
    {
      \begin{cvitems} % Description(s) of tasks/responsibilities
        \item {Développement de systèmes de détection par apprentissage profond pour les images générées par IA, en Python et PyTorch}
        \item {Évaluation systématique des systèmes de reconnaissance faciale face aux attaques deepfake, sur des jeux de données images et vidéos}
        \item {Communication de concepts complexes d'IA et d'apprentissage profond à des publics non techniques (experts forensiques, forces de l'ordre, grand public), en adaptant le message à des profils très variés}
        \item {Encadrement de quatre étudiants en Master sur des projets d'apprentissage automatique, avec accompagnement technique sur l'évaluation de modèles et les bonnes pratiques de développement}
      \end{cvitems}
    }

%---------------------------------------------------------
  \cventry
    {Stagiaire de recherche - Génération de données synthétiques} % Job title
    {Institut de recherche Idiap} % Organization
    {Martigny, Suisse} % Location
    {Juil. 2020 - Déc. 2020} % Date(s)
    {
      \begin{cvitems} % Description(s) of tasks/responsibilities
        \item {Création d'un jeu de données synthétique pour la reconnaissance faciale à l'aide de modèles génératifs, avec mise en place de pipelines de génération et d'automatisation du contrôle qualité}
        \item {Évaluation comparative de plusieurs modèles de reconnaissance faciale sur des données synthétiques et réelles, avec analyse des compromis de performance}
      \end{cvitems}
    }

%---------------------------------------------------------
  \cventry
    {Stagiaire de recherche - Apprentissage automatique audio} % Job title
    {Logitech} % Organization
    {Lausanne, Suisse} % Location
    {Fév. 2019 - Août 2019} % Date(s)
    {
      \begin{cvitems} % Description(s) of tasks/responsibilities
        \item {Développement d'un modèle CNN de classification de scènes audio par apprentissage profond avec TensorFlow pour l'égalisation automatique de haut-parleurs, couvrant le pipeline complet (collecte de données, développement du modèle, évaluation, prototype)}
        \item {Déploiement du modèle sur matériel embarqué, avec optimisation pour le traitement en temps réel en tenant compte des contraintes de latence et de mémoire}
        \item {Présentation des prototypes aux équipes de recherche et d'ingénierie, avec démonstration fonctionnelle illustrant l'intégration du modèle}
      \end{cvitems}
    }

%---------------------------------------------------------
\end{cventries}
